\section*{第 \ref{ch:b3:quantum-statistics} 章——量子统计}

\begin{solution}{exo:b3:quantum-statistics:1}
(a) $\qty{7.1}{eV}$，$\qty{8.2e4}{K}$。(b) $\qty{3.2}{eV}$，
$\qty{3.7e4}{K}$。(c) 取 $m = 3u$：$E_{\text{F}} \approx
\qty{4e-4}{eV}$，$T_{\text{F}} \approx \qty{5}{K}$（理想气体
值；相互作用会使其降低）。(d) 金属电子：在
地球上的任何温度下均简并；$^3$He：仅在几开尔文以下简并；
$T = \qty{300}{K}$ 时远未触及金属的\omterm{thm:b3:quantum-statistics:fermi}{费米能}量尺度。
\end{solution}

\begin{solution}{exo:b3:quantum-statistics:2}
(a) $\int_0^{E_{\text{F}}}\epsilon\,g\,\dd\epsilon/\int g\,\dd
\epsilon = \tfrac35 E_{\text{F}}$，由 $g \propto \sqrt\epsilon$ 可得。
(b) $v_{\text{F}} = \sqrt{2E_{\text{F}}/m_{\text{e}}} =
\qty{1.6e6}{m/s}$；方均根速率 $\sqrt{3/5}\,v_{\text{F}} \approx
\qty{1.2e6}{m/s}$。(c) $\tfrac32 k_{\text{B}}T = \tfrac35
E_{\text{F}}$：$T \approx \qty{3.3e4}{K}$。(d) 无处可去：
所有更低的态均已占满，故既无光子可发射，
也无碰撞能使其减速——有速率而无热量，全因泡利不相容。
\end{solution}

\begin{solution}{exo:b3:quantum-statistics:3}
(a) $P = 0.4 \times \num{8.5e28} \times \qty{7.1}{eV} =
\qty{3.9e10}{Pa} \approx \num{4e5}$ 个大气压。(b) 离子晶格的库仑
吸引作用：金属正是二者之间的
妥协（见\cref{exo:b3:quantum-statistics:11}）。(c) 压缩会使
$E_{\text{F}}$ 陡增（$n^{2/3}$）：金属因而抗压。(d) $K \sim
\tfrac53 P \approx \qty{65}{GPa}$，与实测的
\qty{140}{GPa}相比：铜一半的刚度纯粹源于泡利不相容。
\end{solution}

\begin{solution}{exo:b3:quantum-statistics:4}
(a) $T_{\text{c}} \approx \qty{0.4}{\micro K}$——JILA 的
凝聚体出现于 \qty{170}{nK}：物理本质相同，但
峰值密度略低。(b) $T_{\text{c}} \propto 1/m$：$\times87 \approx
\qty{35}{\micro K}$。(c) 在高密度与如此低温下，
除自旋极化氢外，其余气体均会\emph{凝固}；三体
碰撞会生成分子。关键在于稀薄的亚稳态气体，
其温度足够低，使 $\lambda_T$ 足以抵消稀薄的影响。(d)
$\lambda_{T_{\text{c}}} = h/\sqrt{2\pi mk_{\text{B}}T_{\text{c}}}
\approx \qty{0.3}{\micro m}$：$n\lambda^3 = \num{e20} \times
\num{2.7e-20} \approx 2.7$，与该判据的要求相符。
\end{solution}

\begin{solution}{exo:b3:quantum-statistics:5}
(a) $\gamma T = \beta T^3$ 在 $T^2 = \gamma/\beta$ 时成立；对于铜
$T \approx \qty{3}{K}$。(b) $C_{\text{el}}/C \approx
(\pi^2/2)(300/\num{8.2e4})/3 \approx 0.6\%$。(c) 在 \qty{0.1}{K}
时，晶格项随 $T^3$ 骤降，电子
所携带的热容多出 $\sim10^3$ 倍——量热术变为电子谱学。
(d) $C/T = \gamma + \beta T^2$ 在以
$T^2$ 为横坐标的图中为直线：截距 $\gamma$ 度量\omterm{thm:b3:quantum-statistics:fermi}{费米能}级处的\omterm{prop:b3:schrodinger-three-dimensions:counting}{态密度}，
斜率 $\beta$ 给出德拜温度——同一张图
即可得出每种金属的两个数。
\end{solution}

\begin{solution}{exo:b3:quantum-statistics:6}
(a) $N_{\text{e}} = M/2m_{\text{n}} = \num{6e56}$；$n =
\num{5.5e35}\,\unit{m^{-3}}$；$E_{\text{F}} \approx
\qty{0.24}{MeV}$。(b) 约为电子\omterm{def:b3:relativistic-dynamics:momentum-energy}{静能}的一半（相对论
性边缘），却为 $k_{\text{B}}T$ 的三百倍
（\qty{0.9}{keV}）：既炽热，又处于量子冷态。(c)
$P_{\text{简并}} \approx \qty{9e21}{Pa}$，与 $GM^2/R^4 \sim
\num{e22}\,\unit{Pa}$：量级相同——恒星正是这一
平衡。(d) 质量增大时：引力一侧随 $M^2$ 增长，而
恒星必须收缩（$R \propto M^{-1/3}$），使 $E_{\text{F}}$
趋向相对论性软化。
\end{solution}

\begin{solution}{exo:b3:quantum-statistics:7}
(a) $\rho \approx \qty{5e17}{kg/m^3}$：约为核
密度的两倍——相当于一座城市大小的原子核。(b) $n_{\text{n}}
\approx \qty{3e44}{m^{-3}}$：$E_{\text{F}} \approx \qty{90}{MeV}$，
为 $m_{\text{n}}c^2$ 的十分之一：相对论效应已然显现。(c) 电子的
\omterm{thm:b3:quantum-statistics:fermi}{费米能}超过了 \qty{0.78}{MeV} 的
$\text{e} + \text{p} \to \text{n} + \nu$ 反应的代价：俘获变得
有利可图，恒星以其电子与质子换取了
中子。(d) $\qty{5}{cm^3} \times \qty{5e17}{kg/m^3} =
\qty{2.5e12}{kg}$：一茶匙的质量竟超过一座山脉。
\end{solution}

\begin{solution}{exo:b3:quantum-statistics:8}
(a) 冷却谐振腔只是\emph{移除}光子（腔壁吸收
它们）：由于不存在需要守恒的粒子数，$\mu$ 恒为零，
因而没有多余粒子会堆积进基态模式——气体只是变暗，
而非凝聚。(b) 原子数在每个
实验时间尺度上均守恒。(c) 染料反复吸收并再发射
光子，\emph{而不损失}，在热化
时间尺度上：光子数实际守恒（以及由腔体给出的
有效质量）——凝聚果然出现。
(d) 声子可自由地产生与湮灭，$\mu = 0$：不会
凝聚——晶体中的“声子气”只是冻结掉了。
\end{solution}

\begin{solution}{exo:b3:quantum-statistics:9}
(a) $T_{\text{c}}^{\text{理想}} \approx \qty{3.1}{K}$，与
实测的 $T_\lambda = \qty{2.17}{K}$。(b) 液氦致密且相互作用强，
“理想”模型舍去了会使转变移动的相互作用，
它们还会重塑转变形态。(c) $^3$He 原子是费米子：
无配对则不会凝聚；吸引相互作用
仅在 \qty{2.6}{mK} 以下才束缚成类库珀对，配对的液体
随即发生超流。(d) 九个费米子：$^6$Li 是费米子（而
$^7$Li 是玻色子）——这一对锂同位素使一个实验室可以
在同一炉中研究费米海与凝聚体。
\end{solution}

\begin{solution}{exo:b3:quantum-statistics:10}
(a) 代入 $g \propto \sqrt\epsilon$，并消去 $\beta$。
(b) 积分值给出 $N_{\text{激发}} =
2.612\,V/\lambda_T^3$；令 $N_{\text{激发}} = N$ 定义
$T_{\text{c}}$。(c) 在 $T_{\text{c}}$ 以下，$N_{\text{激发}}(T) =
N(T/T_{\text{c}})^{3/2}$；其余为 $N_0$。(d) 二维情形下
$g(\epsilon)$ 为常数，且 $\int\dd\epsilon/(\eu^{
\beta\epsilon} - 1)$ 在底端发散：激发态
总能容纳全部粒子——基态上不存在宏观
布居（对任意 $T > 0$ 皆然）。
\end{solution}

\begin{solution}{exo:b3:quantum-statistics:11}
(a) $\dd\epsilon/\dd n = \tfrac23 An^{-1/3} - \tfrac13 Bn^{-2/3} =
0$ 处，$n^{1/3} = B/2A$ 为真正的极小值。(b) 取 $A \sim
\hbar^2/m_{\text{e}}$ 与 $B \sim e^2/4\pi\varepsilon_0$，
平衡间距 $n^{-1/3} \sim \hbar^2 4\pi\varepsilon_0/
m_{\text{e}}e^2 = a_0$：金属之所以致密，乃因\omterm{thm:b3:hydrogen-atom:levels}{玻尔半径}
如此。(c) $\epsilon_{\min} \sim -B^2/A \sim -E_{\text{I}}$
——电子伏特量级的结合能，与实测相符。(d) \omterm{thm:b3:quantum-statistics:fermi}{简并压}
是这场妥协中排斥的一半：若无此项，库仑项
将把晶格压缩为一点——物质的体积是泡利的功劳。
\end{solution}

\begin{solution}{exo:b3:quantum-statistics:12}
(a) 活跃电子数为 $g(E_{\text{F}})k_{\text{B}}T$；每个电子的能量：
$\sim k_{\text{B}}T$；求微分得 $E_{\text{额外}} \sim
g k_{\text{B}}^2T^2$。(b) 取 $g(E_{\text{F}}) = 3N/2E_{\text{F}}$
可得 $C \sim 3Nk_{\text{B}}T/T_{\text{F}}$——精确结果中
$\pi^2/2$ 取代了其中的 3。(c) 可翻转自旋数 $\propto T$；每个
自旋的取向贡献 $\propto \mu^2B/k_{\text{B}}T$：$T$ 之因子
相互抵消——泡利磁化率为常数，而居里
磁化率（\cref{exo:b3:canonical-ensemble:11}）随 $1/T$ 下降：一个
简并性的判据。(d) 简并费米子仅在其
表面作出响应：将任一经典结果乘以 $\sim
T/T_{\text{F}}$，或将其限于费米壳层。
\end{solution}

\begin{solution}{pb:b3:quantum-statistics:1}
\textbf{1.} 由分散的物质聚集成恒星所释放的能量
——为负，越致密则越深。
\textbf{2.} $E_{\text{F}} \sim (\hbar^2/m_{\text{e}})
(N_{\text{e}}/R^3)^{2/3}$。
\textbf{3.} $E_{\text{k}} \sim N_{\text{e}}E_{\text{F}} \sim
\hbar^2N_{\text{e}}^{5/3}/m_{\text{e}}R^2$。
\textbf{4.} 极小化 $A/R^2 - B/R$，得 $R = 2A/B$：$R \sim
\hbar^2N_{\text{e}}^{-1/3}/(Gm_{\text{e}}m_{\text{n}}^2\mu_e^2)$。
\textbf{5.} $N_{\text{e}} \propto M$：$R \propto M^{-1/3}$——
增加质量使恒星\emph{收缩}，因而密度与 $E_{\text{F}}$
无限增长：非相对论性的支撑
是靠透支度日。
\textbf{6.} $N_{\text{e}} = \num{6e56}$：$R \sim \qty{2000}{km}$
——与天狼星 B 的 \qty{5800}{km} 量级相当（我们的标度
略去了量级为几的因子）。
\textbf{7.} 当 $E_{\text{F}} \sim m_{\text{e}}c^2$ 时，$n \sim
(m_{\text{e}}c/\hbar)^3 \approx \qty{2e37}{m^{-3}}$，即 $\rho
\sim \num{e10}$--$\num{e11}\,\unit{kg/m^3}$：为
最重\omterm{ex:b3:quantum-statistics:dwarf}{白矮星}的区间。
\textbf{8.} 每个电子：$\epsilon \sim \hbar c\,n^{1/3}$；
总能量：$\hbar cN_{\text{e}}^{4/3}/R$。
\textbf{9.} 两项均 $\propto 1/R$，半径因而消去：
稳定性取决于动能系数是否超过
引力项的系数——这是纯粹的质量判据。
\textbf{10.} 令两者相等，得 $N_{\text{e,临界}} \sim (\hbar c/
Gm_{\text{n}}^2\mu_e^2)^{3/2}$，且 $M_{\text{Ch}} =
\mu_em_{\text{n}}N_{\text{e,临界}}$：即所给公式。
\textbf{11.} $\hbar c/Gm_{\text{n}}^2 = \num{1.7e38}$；
$M_{\text{Ch}} \sim (\num{1.7e38})^{3/2}m_{\text{n}}/4 \approx
\qty{9e29}{kg} \approx 0.5\,M_\odot$——精确计算
将数量级估计修正为 $1.4\,M_\odot$。
\textbf{12.} $\hbar$（压强是量子的），$c$（其
相对论性软化），$G$（对抗者）：一个恒星尺度的数
由三个微观常数构成——量子力学
为含 $10^{57}$ 个粒子的天体立法。
\textbf{13.} 中子简并，半径缩小约
$\sim m_{\text{e}}/m_{\text{n}}$ 倍（另计 $\mu_e$ 的换算因子）：
千米量级——即中子星。
\textbf{14.} 三个相同的常数，处处取 $m_{\text{n}}$，
给出相同的质量标度：中子星的质量上限约为
$2\,M_\odot$；超过此值，没有任何压强定律能够抗衡——形成黑洞。
\textbf{15.} $GM^2/R \approx \qty{5e46}{J}$：数百倍于
太阳一百亿年间的全部输出，却于数秒内释放。
\textbf{16.} 超新星 1987A：二十四个中微子，比光线变亮
早数小时抵达，所携能量恰与预言相符——
坍缩得以直接“听到”。
\textbf{17.} 核心坍缩：大质量恒星的铁核内爆，
越过所有压强下限；外层被吹散。Ia 型：一颗
\omterm{ex:b3:quantum-statistics:dwarf}{白矮星}积累至阈值，点燃其碳，
于热核闪光中将整颗恒星抛散。
\textbf{18.} 触发质量、燃料与能量释放皆相同：
光变曲线几乎一致——经过标定的“炸弹”。
\textbf{19.} 已知绝对亮度，由观测亮度
依平方反比律可得距离；若无
“标准”这一前提，距离阶梯即告崩溃。
\textbf{20.} 亮度过低意味着距离过远：膨胀
正在加速——暗能量，其证据依托于烛光
的一致性（因此须审慎考察其物理过程）。
\textbf{21.} 不相容原理 $\to$ \omterm{thm:b3:quantum-statistics:fermi}{简并压} $\to$ 一个
普适极限质量 $\to$ 标准化的爆发 $\to$
宇宙膨胀的实测加速度。
\textbf{22.} 循数学所指而行：
超过极限的恒星的“荒谬”行为，结果竟是
中子星与黑洞——如今两者均已编目
数以千计。
\textbf{23.} 1915 年，已知物理不允许物质达到每
立方厘米一吨；1926 年，费米–狄拉克统计使之成为
任何冷而致密气体的默认状态：观测结果
有待统计理论，而非相反。
\textbf{24.} \omterm{ex:b3:quantum-statistics:dwarf}{白矮星}能在一系列质量上稳定存在
（$R(M)$ 分支），以及该分支\emph{终止于}
$1.4\,M_\odot$：两者皆绘于该直方图中。
\textbf{25.} $(\hbar c/Gm_{\text{n}}^2)^{3/2}m_{\text{n}} \approx
1.4\,M_\odot$：三个常数决定了最重的\omterm{ex:b3:quantum-statistics:dwarf}{白矮星}——
并由此确定了称量宇宙命运的标准烛光的
触发质量。
\end{solution}
